% ============================================================
% SECTION 1: INTRODUCTION
% ============================================================
% Target: 600--800 words.
%
% Cover:
%   - Why hedging matters (risk transfer, market makers, etc.)
%   - Classical approach: model-based delta hedging and its limitations
%     (model misspecification, discrete rebalancing, transaction costs).
%   - Motivation for a data-driven / model-free approach.
%   - State your research question(s) explicitly.
%   - Brief outline of the paper.
% ============================================================

\section{Introduction}
\label{sec:introduction}

The problem of hedging financial derivatives lies at the heart of modern quantitative
finance. A seller of a derivative security, such as a European call option, faces the
risk that the payoff owed at maturity may exceed the premium received. The classical
approach to managing this risk is \emph{delta hedging}, pioneered by
\citet{BlackScholes1973} and \citet{Merton1973}: one dynamically trades the underlying
asset so that, in continuous time and under the model's assumptions, the hedging error
is identically zero.

In practice, however, continuous trading is impossible, all models are misspecified to
some degree, and every trade incurs a transaction cost. These frictions motivate the
question: \emph{can neural networks learn a near-optimal hedging strategy directly from
simulated (or historical) price data, without explicit knowledge of the underlying
model?}

This paper addresses that question. Specifically, \ldots

% TODO: State research questions explicitly here.

\paragraph{Outline.} The remainder of the paper is organised as follows.
Section~\ref{sec:litreview} reviews the relevant literature.
Section~\ref{sec:methods} describes the models and neural network architecture.
Section~\ref{sec:results} presents and analyses our results.
Section~\ref{sec:conclusion} concludes.
